\section{Introduction}

% Source: .context/chat/Paper Overview (2026-11-19) (2).md, read end-to-end.
% Structure: seven paragraph roles; provisional drafting budget 900--1,100 words.
% Governing revision: in-sample, descriptive and inferential economic measurement.
% Held-out evaluation: trained DNN comparison where needed for generalization.
% Narrative arc: constructed proximity -> frozen measurement instrument ->
% representability/geometry distinction -> spatial economic consequences.

% P1 -- Economic problem: financial proximity as a measurement choice.
% Opening objects: firm characteristics, common-factor exposures, financial risk.
% Motivation: balance sheets, valuation, leverage and business environments.
% Central problem: economically relevant similarity; constructed, unobserved metric.
% Distinctions: common exposure; residual dependence; causal shock transmission.
% Transition: characteristic-based spatial finance as the empirical starting point.

% P2 -- Fernandez as the empirical foundation.
% Anchor: characteristic-history financial distances; nongeographical neighbors.
% Original setting: 126 Latin American firms, 1997--2006; selected financial ratios.
% Replication scope: asset-pricing component; explicit US reimplementation changes.
% Preserved benchmark: original ratios and faithful spatial CAPM specification.
% Extension: Kou--Peng--Zhong S-APT; subsequent heterogeneous/network formulations.
% Claim boundary: empirical spatial association versus conditional pricing theory.
% Identification concerns: kernel choice, normalization, intercept and spatial term.
% Attribution of kernel findings: US reimplementation; original-sample cause unknown.

% P3 -- Modern US empirical foundation and controlled extension.
% Active sources: historical CRSP returns and SEC accounting histories.
% Primary frequency: quarterly fundamentals and Fernandez-style inference.
% Sample description: verified dates, formation rule and realized coverage counts.
% Population accounting: issuer/security distinction; cohort retention through exits.
% Coverage qualifications: linkage, missing fundamentals and terminal returns.
% Contribution: broader market, regimes and characteristic space; beyond row count.
% SPOTTING -- Design-bridge table: Fernandez | US reimplementation | TFM extension.
% Table contents: universe, dates, ratios, distances, weights, models, estimands.
% Transition: a larger dependent panel and the financial estimation burden.

% P4 -- Statistical motivation for frozen representations.
% Illustrative scale only: 1,000 firms x 80 quarters = 80,000 dependent rows.
% Dependence: common shocks, serial persistence, industry and spatial clustering.
% Contrast: financially trained encoder versus imported frozen nonlinear features.
% Conditional advantage: reduced financial-sample representation estimation.
% Full estimation inventory: context, preprocessing, loadings, intercepts,
% factor covariance, innovation variances, spatial coefficients and readouts.
% Native fixed geometry versus supervised calibration: distinct complexity budgets.
% Required qualifications: dependence, identification, selection and uncertainty.
% Statistical-learning claims: conditional; no automatic admissibility from freezing.
% Trained DNN comparison: matched inputs/targets; appropriate held-out evaluation.

% P5 -- Central measurement question: representability versus geometric meaning.
% Economic target: exposures or another explicitly defined economic quantity.
% Representability cue: target constant on encoder fibers; downstream recovery.
% Geometric non-identification cue: information-preserving non-isometric transforms.
% Distinct objects: predictive information; economic recovery; native distances.
% Controlled contrast: conventional ratios | native TFM | calibrated TFM geometry.
% SPOTTING -- Opening proximity exhibit: same firms under ratios and embeddings.
% Exhibit plan: two or three actual firms; neighbors and exposure/ratio comparisons.
% Purpose: visible measurement choice; no assumed advantage for either geometry.
% Formal payoff deferred: representability/geometric non-identification proposition.

% P6 -- Spatial economic consequences and interpretability.
% Economic chain: geometry -> weights -> spatial model -> covariance/pricing/risk.
% Model distinctions: SAR transformed exposures; SEM residual covariance;
% factor-based SLX modified exposures; specification-specific pricing restrictions.
% Separate hypotheses: common-exposure geometry and residual-dependence structure.
% Theory cue: actual graph normalization, multiplier stability, diversification.
% Maintained economic assumptions: factor structure, trading and no asymptotic arbitrage.
% Inference endpoints: spatial coefficients, pricing restrictions, visible intervals.
% Risk endpoints: systematic/residual covariance; portfolio variance; model-implied VaR.
% Risk scope: in-sample structure and historical fit; explicit distribution assumptions.
% Interpretation: PCoA directions/subspaces, ratio/exposure alignment and stability.
% Further tests: factor-mimicking construction, spanning/pricing and TFM attribution.
% Claim boundaries: coordinates versus priced factors; attribution versus mechanism;
% covariance fit/stability versus no-arbitrage pricing restrictions.

% P7 -- Four contributions, evidence-led findings and roadmap.
% C1: modern Fernandez foundation; characteristic-space comparison and S-APT extension.
% C2: frozen versus financially trained representations; conditional complexity argument.
% C3: structured spatial inference, covariance and portfolio risk; PCoA and explanation.
% C4: economic/representation-learning bridge; AI as a measurement instrument.
% Findings slots: replication; frozen/trained comparison; native/calibrated geometry;
% pricing/covariance/risk effects with intervals; economically interpretable structure.
% Evidence labels: supported | unsupported | unidentified; negative findings retained.
% Findings source: corrected active-study outputs; no inherited withdrawn-study numbers.
% Roadmap: literature review -> methodology -> empirical design and data ->
% results -> conclusion; technical proofs and extended diagnostics in appendices.
